\documentclass[11pt]{article}
\usepackage[a4paper,margin=1in]{geometry}
\usepackage{fontspec}
\usepackage[boldfont=false]{xeCJK}
\setCJKmainfont{FandolSong-Regular.otf}
\usepackage{amsmath}
\usepackage{amssymb}
\usepackage{array}
\usepackage{booktabs}
\usepackage{graphicx}
\usepackage{adjustbox}
\usepackage{enumitem}
\setlist[itemize]{topsep=2pt,partopsep=0pt,parsep=0pt,itemsep=2pt,leftmargin=1.4em}
\setlist[enumerate]{topsep=2pt,partopsep=0pt,parsep=0pt,itemsep=2pt,leftmargin=1.6em}
\usepackage{parskip}
\usepackage{fvextra}
\fvset{breaklines=true,breakanywhere=true,fontsize=\small}
\setlength{\parindent}{0pt}

\begin{document}

\begin{center}
  {\LARGE\bfseries Databases}\\[0.35em]
  {\large A Level Computer Science}
\end{center}
\vspace{0.6em}

\section*{File-based storage and its limits}

Before databases, programs stored data in \textbf{flat files} 平面文件 — usually one file per program. This is fine for small data but breaks down at scale.

\subsection*{Limitations}

\begin{itemize}
\item \textbf{data redundancy} 数据冗余 — the same data (a customer's address) is held in several files.

\item \textbf{data inconsistency} 数据不一致 — redundant copies updated separately get out of sync.

\item \textbf{data dependence} — programs are tied to the file format; change the format and every program must be rewritten.

\item hard to enforce \textbf{integrity} 完整性, hard to share safely, weak querying, and weak per-field security.

\end{itemize}

\par\medskip
\noindent\begin{minipage}{\linewidth}
\centering
\adjustbox{max width=\linewidth, max totalheight=0.42\textheight}{%
  \includegraphics{/Users/joschka/Desktop/Programming/eXercise/output/eXam/handouts/a_level_computer_science/assets/08-p214-file-based-approach.pdf}}\\[0.3em]
\emph{The file-based approach: each program keeps its own files}
\end{minipage}\par\medskip

A \textbf{relational database} 关系数据库 fixes these by storing data in \textbf{tables} managed by one piece of software (the DBMS) that all programs use.

\par\medskip
\noindent\begin{minipage}{\linewidth}
\centering
\adjustbox{max width=\linewidth, max totalheight=0.42\textheight}{%
  \includegraphics{/Users/joschka/Desktop/Programming/eXercise/output/eXam/handouts/a_level_computer_science/assets/08-p215-database-approach.pdf}}\\[0.3em]
\emph{The database approach: one DBMS serves all the programs}
\end{minipage}\par\medskip

\section*{Relational model — terms}

\begin{itemize}
\item \textbf{table} 表 (relation) — a grid of rows and columns; one table per type of \textbf{entity} 实体 (e.g. \texttt{CUSTOMER}).

\item \textbf{record} 记录 (row) — one row; one instance of the entity.

\item \textbf{field} 字段 (column) — one column; one piece of information about each record.

\item \textbf{primary key} 主键 — a field (or fields) that \textbf{uniquely identifies} each record; never null or duplicated.

\item \textbf{foreign key} 外键 — a field whose value matches the primary key of another table, linking the two.

\item \textbf{composite key} 复合键 — a primary key made of two or more fields together.

\item \textbf{candidate key} 候选键 — any field(s) that could be the primary key.

\item \textbf{referential integrity} 参照完整性 — every foreign-key value must match an existing primary key (no orphan records).

\end{itemize}

A table is written in shorthand with the primary key underlined and foreign keys noted:

\begin{Verbatim}
CUSTOMER(CustomerID, Name, Phone)
ORDER(OrderID, CustomerID, OrderDate)   -- CustomerID is FK → CUSTOMER
\end{Verbatim}

\section*{Entity-relationship (E-R) diagrams}

An \textbf{entity-relationship diagram} 实体关系图 shows the structure: each entity is a rectangle, each relationship a line, with the \textbf{cardinality} 基数 marked at each end:

\begin{itemize}
\item \textbf{one-to-one} (1:1).

\item \textbf{one-to-many} 一对多 (1:M) — each Customer has many Orders; each Order has one Customer.

\item \textbf{many-to-many} (M:N) — Students take many Courses, and Courses have many Students.

\end{itemize}

\par\medskip
\noindent\begin{minipage}{\linewidth}
\centering
\adjustbox{max width=\linewidth, max totalheight=0.42\textheight}{%
  \includegraphics{/Users/joschka/Desktop/Programming/eXercise/output/eXam/handouts/a_level_computer_science/assets/08-p218-er-diagram.pdf}}\\[0.3em]
\emph{An E-R diagram: one class has many students}
\end{minipage}\par\medskip

\par\medskip
\noindent\begin{minipage}{\linewidth}
\centering
\adjustbox{max width=\linewidth, max totalheight=0.42\textheight}{%
  \includegraphics{/Users/joschka/Desktop/Programming/eXercise/output/eXam/handouts/a_level_computer_science/assets/08-p218-cardinality-notation.pdf}}\\[0.3em]
\emph{Crow's-foot symbols for the cardinality of a relationship}
\end{minipage}\par\medskip

A many-to-many relationship cannot be stored directly. Break it into two one-to-many relationships through a \textbf{link table} 连接表 holding the two foreign keys:

\begin{Verbatim}
ENROLMENT(StudentID, CourseID, EnrolmentDate)
\end{Verbatim}

\section*{Normalisation}

\textbf{Normalisation} 规范化 organises tables to \textbf{cut redundancy and inconsistency}, going through \textbf{normal forms} 范式 in order.

\begin{itemize}
\item \textbf{First normal form (1NF)} — every field holds a single (\textbf{atomic} 原子) value, with no repeating groups, and a primary key.

\item \textbf{Second normal form (2NF)} — in 1NF, and every non-key field depends on the \textbf{whole} primary key (only matters for a composite key).

\item \textbf{Third normal form (3NF)} — in 2NF, and every non-key field depends \textbf{only} on the primary key, not on another non-key field (no \textbf{transitive dependency} 传递依赖).

\end{itemize}

A 3NF design stores each fact once, so insert/update/delete anomalies disappear. The trade-off is more tables and more joins. Aim for 3NF.

To produce a 3NF design: find the entities and their attributes; choose a primary key for each; split repeating/non-atomic fields (1NF); split fields depending on part of a composite key (2NF); split fields depending transitively on the key (3NF); add foreign keys for the relationships.

\section*{Database Management System (DBMS)}

A \textbf{DBMS} 数据库管理系统 manages the database centrally. Features that fix the file-based limits:

\begin{itemize}
\item \textbf{data dictionary} 数据字典 — a description of every table, field, type and key; programs query it instead of hard-coding the structure.

\item redundancy/consistency control — each fact stored once.

\item \textbf{concurrent access} 并发访问 control — locks and transactions let many users work at once.

\item \textbf{backup} 备份 and recovery; security and per-user permissions.

\item \textbf{integrity rules} — keys, unique and range constraints, enforced centrally.

\item \textbf{transactions} 事务 — a group of operations that all succeed or all fail.

\item \textbf{views} 视图 — virtual tables that show each user "their" slice of the data.

\end{itemize}

Its tools include a data-dictionary editor, a query builder, a forms builder, a report generator, user management, and an SQL editor.

\section*{DDL and DML}

\textbf{SQL} 结构化查询语言 (Structured Query Language) has two halves:

\begin{itemize}
\item \textbf{Data Definition Language} 数据定义语言 (DDL) — creates or changes the \textbf{structure} (tables, keys, constraints).

\item \textbf{Data Manipulation Language} 数据操纵语言 (DML) — works with the \textbf{data} (insert, update, delete, \textbf{query} 查询).

\end{itemize}

\subsection*{DDL basics}

\begin{Verbatim}
CREATE TABLE CUSTOMER (
  CustomerID INTEGER PRIMARY KEY,
  Name VARCHAR(50) NOT NULL,
  Phone VARCHAR(20)
);
\end{Verbatim}

Add a foreign key:

\begin{Verbatim}
CREATE TABLE ORDER (
  OrderID INTEGER PRIMARY KEY,
  CustomerID INTEGER,
  OrderDate DATE,
  FOREIGN KEY (CustomerID) REFERENCES CUSTOMER(CustomerID)
);
\end{Verbatim}

Modify and drop:

\begin{Verbatim}
ALTER TABLE CUSTOMER ADD Email VARCHAR(100);
DROP TABLE CUSTOMER;
\end{Verbatim}

Common types: \texttt{INTEGER}, \texttt{VARCHAR(n)}, \texttt{CHAR(n)}, \texttt{DATE}, \texttt{BOOLEAN}, \texttt{DECIMAL(p, s)}.

\subsection*{DML basics}

Query with \texttt{SELECT}:

\begin{Verbatim}
SELECT Name, Phone
FROM CUSTOMER
WHERE City = 'London'
ORDER BY Name ASC;
\end{Verbatim}

\texttt{SELECT} lists fields, \texttt{FROM} names the table, \texttt{WHERE} filters rows, \texttt{ORDER BY} sorts.

A \textbf{join} 连接 combines two tables using a foreign-key relationship:

\begin{Verbatim}
SELECT C.Name, O.OrderDate
FROM CUSTOMER C INNER JOIN ORDER O
  ON C.CustomerID = O.CustomerID
WHERE O.OrderDate >= '2024-01-01';
\end{Verbatim}

\textbf{Aggregate functions} 聚合函数 (\texttt{COUNT}, \texttt{SUM}, \texttt{AVG}, \texttt{MIN}, \texttt{MAX}) are often used with \texttt{GROUP BY}:

\begin{Verbatim}
SELECT CustomerID, COUNT(*) AS NumOrders
FROM ORDER
GROUP BY CustomerID;
\end{Verbatim}

Insert, update, delete:

\begin{Verbatim}
INSERT INTO CUSTOMER (CustomerID, Name, Phone)
VALUES (101, 'Ada Lovelace', '020-1234-5678');

UPDATE CUSTOMER SET Phone = '020-9999-0000' WHERE CustomerID = 101;

DELETE FROM CUSTOMER WHERE CustomerID = 101;
\end{Verbatim}

Always put a \texttt{WHERE} clause on \texttt{UPDATE} and \texttt{DELETE}, or the change hits \textbf{every} row.

\subsection*{Tips for exam SQL}

\begin{itemize}
\item use the exact table and field names from the question.

\item quote strings with single quotes (\texttt{'Smith'}); don't quote numbers.

\item comparisons: \texttt{=}, \texttt{\textless{}}, \texttt{\textgreater{}}, \texttt{\textless{}=}, \texttt{\textgreater{}=}, \texttt{\textless{}\textgreater{}}.

\item \texttt{LIKE 'A\%'} matches anything starting with A (\texttt{\%} = any string, \texttt{\_} = one character); \texttt{IN (1,2,3)}; \texttt{BETWEEN 10 AND 20}.

\item combine conditions with \texttt{AND} / \texttt{OR} / \texttt{NOT}, and end each statement with a semicolon.

\end{itemize}


\end{document}
